\section{\texorpdfstring{Innovativeness and originality}{Innovativeness and originality }}\label{innovativeness-and-originality}

\subsection{Context}\label{context-1}

Cognitive robotics has the potential to revolutionize manufacturing by
enabling robots to not only execute tasks but also understand and adapt
to complex, dynamic environments.

With cognitive capabilities, robots could autonomously make decisions,
collaborate more effectively with human workers, and handle unexpected
situations, significantly boosting manufacturing productivity,
efficiency, and flexibility. For example, a cognitive robot could adjust
its behavior in real time based on new data, optimize processes on the
fly, or even learn new tasks without extensive reprogramming.
\textbf{However, several hardware and software tools are still missing
to realize this vision fully.}

On the hardware side, more powerful yet efficient processors are needed
to handle the computational demands of real-time decision-making and
learning, along with advanced sensors for better environmental
perception and interaction. On the software side, there is a need for
more sophisticated frameworks that integrate machine learning, natural
language understanding, and advanced planning algorithms. Furthermore,
developing scalable and flexible software architectures that allow easy
reconfiguration and learning across diverse tasks remains a significant
challenge in the field.

The disruptive innovation for \textbf{Plan4AIR} consists of the
\textbf{systematic approach to cognitive robotics} to realize the first
cognitive industrial robot able to interact with humans and solve every
task autonomously based on chatting with the process expert.

\subsection{Artificial Robotic Intelligence: the innovation and the
rocket
science}\label{artificial-robotic-intelligence-the-innovation-and-the-rocket-science}

According to the applied nature of the funding schema, \brand{Plan4ARI}
aims first to deploy a cognitive robot that exploits the best from the
fundamental research. Generative AI, deliberative AI, deep learning
semantic sensing, machine learning for motion planning, and neural
control are recent scientific breakthroughs in each research area.

No research has tried to combine the best abilities to allow the robot
to achieve a step-change in its cognitive ability. Robots are not only
AI flesh out but also complex systems; the balance between the different
modules is the real key to success.

To do so, \brand{Plan4ARI} aims at a \textbf{pragmatic} approach:

\begin{itemize}
\item
  We do not reinvent the wheel. We want to exploit the best solutions
  available in the field -- and there is a huge quantity of high-quality
  tools in academia (open-source community) and off-the-shelves tools
\item
  We exploit the MI.RA Dexter COMAU platform, which allows low-frequency
  async communication with the robot to move it and access the
  industrial automation connected to the robot. This will make the
  development of the modules for natural language programming and
  adaptive motion planning straightforward.

  \begin{itemize}
  \item
    \textbf{NOTE: all the scientific outcomes will be independent of the
    COMAU robot controller and extendable to all industrial
    manipulators. The adoption from COMAU of the Plan4ARI platform will
    grant an essential advance to the Italian sector worldwide, allowing
    the Italian industrial robotic industry to be a breakthrough in the
    market}
  \end{itemize}
\item
  We will perform \textbf{fundamental research} only where there is a
  need. We will innovate in the field of

  \begin{itemize}
  \item
    generative AI for the Knowledge Base foundation model and the PDDL
    file.
  \item
    sensor fusion of radar sensing with 3D camera systems (a completely
    new research field).
  \end{itemize}
\item
  We will perform \textbf{industrial research} to deploy the best and
  most competitive tools for cognitive robotics. We will be disruptive
  in the development of robust and high-quality methodologies in the
  field of:

  \begin{itemize}
  \item
    deliberative AI tools integration with STNU
  \item
    adoption of MCTS to optimize the task plan using virtual execution
    of the tasks.
  \item
    sample-based informed set motion planner to allow runtime motion
    planning
  \item
    fusion of sample-based techniques with high-performing MPC to grant
    robustness in the task execution and exploitation of residual
    degrees of freedom of the setup
  \item
    Deploy the first industrial controller, exploiting neural control
    design and prescribed performance constraint control to grant
    robustness and safety in all working conditions.
  \end{itemize}
\end{itemize}

It is worth noting that adopting AI techniques in robotics is now at a
critical stage. The innovation of ideas claimed in this project could be
not so disruptive in 6 months, i.e., at the project start. The
generative AI tools are evolving with impressive velocity, and an update
of the scientific plan will be necessary once the project starts.

\subsection{Artificial Robotic Intelligence: a framework to cognitive
robotics}\label{artificial-robotic-intelligence-a-framework-to-cognitive-robotics}

One crucial innovation presented by \textbf{Plan4AIR} demonstrates that
the insertion of the proper AI tools at the correct robot control level
is the key to having a completely autonomous mobile manipulator able to
solve the problem when asked to.

The project\textquotesingle s goal is that the motion control framework
is as important as the single methodology. The full-working system will
solve complex applications in a considerably short time.

The project will demonstrate that mobile industrial robots can be real
companions of human beings. The trustworthiness investigation will be
crucial at this point. These investigations are always limited to
laboratory setups and small payload robots. Only a few works deal with
high-payload industrial robots\footnote{M. Rehm, I. Pontikis and K.
  Hald, "Automatic Trust Estimation From Movement Data in Industrial
  Human-Robot Collaboration Based on Deep Learning," 2024 IEEE
  International Conference on Robotics and Automation (ICRA), Yokohama,
  Japan, 2024, pp. 11245-11251, doi: 10.1109/ICRA57147.2024.10610822.},
but with a very narrow application. The role of COMAU in this evaluation
will allow us to achieve an extraordinary impact in such a research
field. Accessing the experience of one of the largest and most
innovative industrial robot producers worldwide and connecting with its
vast network of users will substantiate the research dramatically.

\begin{itemize}
\item
  The psychological approach to the evaluation of the interaction will
  drive the deployment.
\end{itemize}

The work package WP6 is designed to evaluate the interaction and extract
the critical aspects of the cognitive processes that human beings
activate when collaborating with the robot.

This cognitive evaluation of the interaction is the counterpart of
having cognitive robots, and it is a powerful tool for designing the
task and motion planner.

The definition of the paths, the sequence of the task, and the modality
followed by the robot when it must slow down to preserve safety or hand
something to the operators are fundamental factors for the definition of
awareness in human operators.

The control framework will integrate the outcome of the trustworthiness
investigation.

The investigation will allow the deployment of task and motion planning
rules to improve users\textquotesingle{} acceptance of such
technologies.